% SPDX-License-Identifier: Apache-2.0
% covers: Station2, Station3, Station4, Station5, Station6, Station7, Station8, Station9, Station10
\datasheet{Station}{Reservation stations of two to ten inputs: operands gathered by tag, sent as one line.}

\dsfacts{\code{lib/parts/src/station.rs}; generator \code{station\_text} in \code{lib/macros/lib.rs}}%
{\code{txhdl\_parts::station::Station2} to \code{Station10}}%
{\code{//docs:station}, \code{//docs:runtime}, \code{//docs:examples}}

\subsection*{Function}

A reservation station gathers the operands of an operation that arrive
on separate channels, out of order, and fires when all of them have.
Each input is a channel of \code{Tagged<TB, T>}: a tag of \code{TB}
bits above a value. The station has a line per tag value, with a cell
per input. An input lands in its cell of the line its tag names.

A line whose every cell holds is sent on the one output as a
\code{Line}\textit{N}: the tag, and a value per input, \code{v0}
upwards. \code{station!(StationN,
N)} writes each station with its line struct beside it, since the
lowering reads a body and not a loop over inputs. \code{station.rs}
writes \code{Station2} to \code{Station10}; the macro itself accepts 2
to 16.

\subsection*{Parameters}

\begin{center}\footnotesize
\begin{tabular}{@{}lp{0.7\textwidth}@{}}
\toprule
Parameter & Meaning \\
\midrule
\code{TB} & The tag width, in bits. It is two at least. \\
\code{L} & The count of lines, which is \code{1 << TB}. \\
\code{T}\textit{i} & The value type of input \textit{i}, for \textit{i}
from 0 to \textit{N} less one. The types may differ. \\
\bottomrule
\end{tabular}
\end{center}

\code{L} is stated rather than computed, because a computed width in a
type needs nightly Rust. The tag has two bits at least, since a one-bit
value is a single wire in VHDL and cannot index a line. The tables
below use \code{Station3<4, 16, U<2>, U<16>, U<16>>}, the station of
\code{ex\_compose}: sixteen lines, an operation and two operands.

\dstables{Station}

\subsection*{Behaviour}

Each input \textit{i} has \code{mem}\textit{i}, its cells, a word per
line, and \code{occ}\textit{i}, a bit per line that says which cells
hold. Each cycle:

\begin{itemize}
\item An input's cell is free when its tag's bit of the input's mask
is clear.
\item Taking an input completes its line when, for every other input,
that line's cell holds, or that input offers the same tag now and its
own cell is free.
\item The line sent is the completing input's of lowest index, if any
input completes and the output has room. At most one line leaves per
cycle.
\item An input is taken when its cell is free and taking it completes
nothing, or completes the line being sent. An input's \code{ready} is
that take.
\item So an input never waits on another input's cell. It waits on its
own cell, on the output's room, or while an input of lower index
completes a different line in the same cycle.
\item A taken input writes its cell and sets its bit. The line sent
clears its bit in every mask.
\item The values sent come from the cells. The value of an input that
completes the line in this cycle goes straight through instead.
\end{itemize}

\subsection*{Verification}

\begin{itemize}
\item Two tests in \code{txhdl\_parts}, in
\code{//lib/parts:parts\_test}. \code{station2\_follows\_the\_rule} runs
\code{Station2<2, 4, U<8>, U<16>>} against the rule written with loops
for 400 cycles of random offers and takes. It checks both masks after
every cycle and every line the sink takes.
\code{station3\_lowest\_input\_wins} sets up two lines that complete in
one cycle, and checks that the lower input's line goes first.
\item \code{ex\_station} runs a \code{Station2<2, 4, U<8>, U<16>>}
beside \code{station.v}, the same station written by hand in Verilog,
and compares every line the two send. The build simulates the
station's netlist against that run under nvc and Verilator:
\code{//docs:station\_sim\_station\_tb\_test} and
\code{//docs:station\_vsim\_test}.
\item \code{ex\_compose} lowers \code{Top}, which holds the
\code{Station3} of the tables and the stream ALU as children. The build
simulates the netlist of the top against that run:
\code{//docs:compose\_sim\_top\_tb\_test} and
\code{//docs:compose\_vsim\_test}.
\end{itemize}

\subsection*{Used by}

\code{Station2} in the example \code{ex\_station}. \code{Station3} in
the example \code{ex\_compose}, in front of the stream ALU, whose
requests are \code{Line3} values.

\subsection*{Limits}

A station is not an arbiter. If an input's next offer names a cell it
already holds, it waits until that line completes. If the other inputs
never offer that tag, it waits for ever, as \code{//docs:station}
states. Each input must, in time, see every tag its lines wait on. One
line leaves per cycle at most, and a line waits for room on the output;
\code{//docs:station} puts a \code{Fifo} behind the station where the
consumer takes in bursts.
